\documentclass[14pt]{extarticle}
\usepackage[utf8]{inputenc}
\usepackage{amsmath, amssymb}
\usepackage{indentfirst}
\usepackage{url}
\usepackage{graphicx}
\usepackage{geometry}
\usepackage{csquotes}
\usepackage{longtable}
\usepackage{booktabs}
\usepackage{array}
\usepackage{tikz}
\usepackage{graphviz}
\usetikzlibrary{positioning}
\usepackage{colortbl}
\newcommand{\colormidrule}{%
  \arrayrulecolor{black!30}\midrule\arrayrulecolor{black}%
}
\usepackage{caption}
\usepackage{tabularx}
\usepackage{ragged2e}
\usepackage{float}
\usepackage[hidelinks]{hyperref}
\hypersetup{
    colorlinks=true,
    citecolor=blue,
    linkcolor=black,
    filecolor=magenta,      
    urlcolor=blue,
    pdftitle={investigation-filesytem},
    pdfpagemode=FullScreen,
}
\usepackage{xurl}
\usepackage{fontspec}
\usepackage{etoolbox}
\usepackage{fancyhdr}
\newcommand{\diagramlabel}{}
\fancypagestyle{page:diagram:continued}{%
  \fancyhf{}%
  \fancyfoot[L]{\diagramlabel}
  \fancyfoot[C]{\thepage}%
  \fancyfoot[R]{Continued on next page. \bfseries Turn over.}%
  \renewcommand{\headrulewidth}{0pt}%
  \renewcommand{\footrulewidth}{0pt}%
}

\fancypagestyle{page:diagram:last}{%
  \fancyhf{}%
  \fancyfoot[L]{\diagramlabel}
  \fancyfoot[C]{\thepage}%
  \fancyfoot[R]{\bfseries End of diagram.}%   % or leave the [R] slot empty
  \renewcommand{\headrulewidth}{0pt}%
  \renewcommand{\footrulewidth}{0pt}%
}
\usepackage{pgfgantt}
\setganttlinklabel{s-s}{Start to start}
\setganttlinklabel{f-s}{Finish to start}
\setganttlinklabel{f-f}{Finish to finish}
\usepackage{pdfpages}
\setmainfont{Times New Roman}
% \renewcommand{\familydefault}{\sfdefault}
\setsansfont{CMU Bright}[
  BoldFont={CMU Bright SemiBold},
  ItalicFont={CMU Bright Oblique},
]
\usepackage[russian]{babel}         
\usepackage[backend=biber, style=gost-numeric]{biblatex}
\addbibresource{references.bib}
\DeclareLanguageMapping{russian}{russian-gost}

\usepackage{titlesec}
\titleformat{\section}
  {\large\bfseries}
  {\scshape{\thesection}. }
  {0pt}
  {\scshape}

\titleformat{\subsection}
  {\large\bfseries}
  {\scshape{\thesubsection}. }
  {0pt}
  {\scshape}


\titleformat{\subsubsection}
  {\bfseries}
  {{\thesubsubsection}. }
  {0pt}
  {}

\titlespacing*{\section}{0pt}{2ex plus 1ex minus .2ex}{2ex plus .1ex}
\titlespacing*{\subsection}{0pt}{2ex plus 1ex minus .2ex}{2ex plus .1ex}
\titlespacing*{\subsubsection}{0pt}{2ex plus 1ex minus .2ex}{0.75ex plus .2ex}
\titlespacing*{\paragraph}{0pt}{0.25ex plus 0.25ex minus .2ex}{0.5em}

\usepackage{xcolor}
\definecolor{phasegray}{HTML}{555D63}
\definecolor{taskgreen}{HTML}{19C934}
\definecolor{taskblue}{HTML}{0067C8}
\definecolor{taskpurple}{HTML}{A600D7}
\definecolor{taskorange}{HTML}{FF8500}
\definecolor{taskpink}{HTML}{F5005A}
\definecolor{gridgray}{HTML}{D8DDE1}
\tikzset{ganttgrid/.style={draw=gridgray,line width=0.35pt}}
\setganttlinklabel{f-s}{}

\usepackage{minted}
\setminted{
    style=xcode,
    fontsize=\footnotesize,
    linenos=true,
    breaklines=true,
}

\renewcommand{\theFancyVerbLine}{%
    \footnotesize
    \textcolor{black}{\arabic{FancyVerbLine}:}
}


\graphicspath{{./images/}}
\geometry{a4paper, margin=1in}
\captionsetup{
    font=small,
    labelfont=bf,
    justification=centering,
    margin=2cm,
    labelsep=endash
}

% Пропуски перед параграфами

\setlength{\parskip}{0em}
\setlength{\parindent}{2em}

\makeatletter
\renewcommand{\verbatim@font}{\ttfamily\normalsize}
\makeatother
\let\oldtexttt\texttt
\renewcommand{\texttt}[1]{{\normalsize\oldtexttt{#1}}}

\urlstyle{same}

\setcounter{secnumdepth}{0}

% Форматы таблиц

\makeatletter

\newif\if@tablenumpending

\newcommand{\tablenumber}[1]{%
    \gdef\@pendingtabnum{Таблица~#1}%
    \@tablenumpendingtrue
}

\newcommand{\twelvept}{
  \fontsize{12pt}{14.4pt}\selectfont
}

\newcommand{\tableheading}[1]{%
    \if@tablenumpending
        \begingroup
        \captionsetup{width=\textwidth, margin=0pt}
        \caption*{%
            \makebox[\textwidth][r]{\@pendingtabnum}\\
            \makebox[\textwidth][c]{\textbf{#1}}%
        }%
        \endgroup
        \@tablenumpendingfalse
    \else
        \caption*{\textbf{#1}}%
    \fi
}

\makeatother

\begin{document}

\pagenumbering{gobble}

\begin{titlepage}
  \small УДК: \ldots \normalsize
  \begin{center}
    \uppercase{\textbf{Производство высокопроизводительных веб-решений на языке C++}}
    \vspace{0.5cm}
    
    \textbf{Н. А. Сугоняко}, \textit{магистрант группы ИВТ/м-26-1-о} \\
    \textbf{К. В. Сергеев}, \textit{магистрант группы ИВТ/м-26-1-о} \\
    Федеральное государственное автономное образовательное учреждение высшего образования "<Севастопольский государственный университет (СевГУ)">, Севастополь, Россия \\
    \vspace{0.25em}
    \textbf{Г. А. Сафьянников}, \textit{студент бакалавриата} \\
    Международный Колледж Импингтон, Кембридж, Великобритания
  \end{center}
  \vspace{-1.75em}
  %
  % Аннотация и ключевые слова

  {\twelvept{}
    \begin{center}
      \textit{\paragraph{\twelvept{}\textit{Аннотация.}} щавель.}
    \end{center}
    \vspace{-2em}
    \begin{center}
      \textit{\textbf{Ключевые слова:} щавель.}
    \end{center}
  }
  
  %Английский

  \vspace{0em}

  \begin{center}
    \uppercase{\textbf{Producing high-performance web-solutions in C++}}
    \vspace{0.5cm}
    
    \textbf{N. A. Sugonyako}, \textit{magistrate student in group IVT/m-26-1-o} \\
    \textbf{K. V. Sergeev}, \textit{magistrate student in group IVT/m-26-1-o} \\
    Federal Autonomous State Educational Institution of Higher Education ``Sevastopol State University'', Sevastopol, Russia \\
    \vspace{0.25em}
    \textbf{H. O. Safiannikov}, \textit{prospective student} \\
    Impington International College, Cambridge, United Kingdom
  \end{center}

  {\twelvept{}
    \begin{center}
      \textit{\paragraph{\twelvept{}\textit{Abstract.}} щавель.}
    \end{center}
    \vspace{-2em}
    \begin{center}
      \textit{\textbf{Keywords:} brazier, C++, web-development, HTTP, performance, coloured output.}
    \end{center}
  }
\end{titlepage}

\phantomsection

\pagebreak

\pagenumbering{arabic}
\setcounter{page}{1}

\section{Введение}\label{sec:intro}

Современное развитие информационных технологий сопровождается постоянным увеличением количества веб-сервисов и объёмов обрабатываемых ими данных. Широкое развитие искусственного интеллекта в недавнее время заострило нужду в надёжной и высокопроизводительной серверной инфраструктуре. Конкретно, ИИ-агенты, способные создавать сотни запросов одновременно, серьёзно повышают нагруженность настоящей инфраструктуры, и представляют собой новое испытание в оптимизации. Современные веб-приложения должны обеспечивать обработку значительного количества запросов, поддерживать большое число одновременно подключённых клиентов и сохранять стабильное время отклика при изменении интенсивности нагрузки.

Одной из фундаментальных проблем серверной разработки является обеспечение масштабируемости при ограниченных вычислительных ресурсах. Она получила известность в связи с формированием концепции C10K, рассматривающей возможность эффективного обслуживания десяти тысяч одновременных подключений на одном сервере~\autocite{dong:c10k-problem_2009}. Дальнейшее развитие вычислительных технологий привело к появлению задачи C10M, связанной с обслуживанием миллионов одновременных соединений~\autocite{song:c10m-problem_2023}. Обе концепции демонстрируют необходимость совершенствования механизмов сетевого взаимодействия, управления памятью и распределения вычислительной нагрузки между процессорными ядрами.

При этом производительность серверных приложений не может быть определена исключительно количеством обрабатываемых запросов в секунду. Существенное значение имеют задержка ответа, устойчивость к высокой конкурентной нагрузке, потребление оперативной памяти и процессорного времени, а также эффективность взаимодействия с внешними системами, включая базы данных и сетевые службы. Следовательно, разработка высокопроизводительных серверных приложений требует комплексного подхода, учитывающего как особенности языка программирования, так и архитектуру такого приложения.

Одним из языков, представляющих широкие возможности для создания производительных серверных систем, является C++. Его особенности, включая компиляцию в машинный код, управлением временем жизни объектов, низкоуровневый доступ к системным ресурсам и современные средства конкурентного программирования, позволяют разрабатывать приложения с относительно небольшими накладными расходами. Появление корутин в стандарте C++20 также упростило реализацию асинхронных вычислений, позволяя разделять операции ожидания ввода-вывода и непосредственную обработку запросов. Вместе с тем использование C++ предполагает внимательное отношение к управлению ресурсами, что делает выбор архитектуры и используемых библиотек особенно значимым.

Несмотря на наличие множества языков и платформ для серверной разработки, экосистема C++ располагает рядом полнофункциональных фреймворков и низкоуровневых сетевых библиотек, ориентированных на создание высокопроизводительных приложений. Среди них можно выделить Drogon, Oat++, userver, Crow и brazier, предлагающие различные подходы к организации асинхронного ввода-вывода, маршрутизации запросов и взаимодействию с базами данных.

Целью настоящей работы является рассмотрение архитектурных и технологических принципов разработки высокопроизводительных веб-решений на языке C++ c анализом существующих инструментов и практической реализацией данных принципов на примере кроссплатформенного фреймворка brazier. Для её достижения рассматриваются основные факторы, влияющие на производительность веб-приложений, анализируются возможности языке C++ в контексте серверной разработки, проводится сравнительный обзор существующих C++ фреймворков и исследуются архитектурные особенности brazier. Отдельное внимание уделяется асинхронной обработке сетевых операций, организации многопоточного выполнение, взаимодействию с базами данных и механизмами управления сетевыми соединениями.

% КОЛЯ vvv (слоп граната)

%Разработка высокопроизводительных веб-решений на C++ связана с необходимостью одновременно обеспечить кроссплатформенность, предсказуемость задержек и низкое потребление ресурсов. Проблема C10K/C10M показывает, что удержание большого числа одновременных соединений требует эффективной событийной модели и экономного расхода памяти на каждое соединение~\autocites{dong:c10k-problem_2009}{song:c10m-problem_2023}. Управляемые среды добавляют недетерминированные паузы сборки мусора, а интерпретируемые — накладные расходы на выполнение. Раздутый объём занимаемой памяти в облачной среде напрямую увеличивает стоимость vCPU и гигабайт-часов. Поэтому представляет интерес анализ инструментов и подходов, которые позволяют разрабатывать такие решения на C++ без избыточных затрат.
%
%Одним из ключевых элементов такого стека являются библиотеки Boost. Они предоставляют кроссплатформенные компоненты для асинхронного ввода-вывода (Boost.Asio), работы с HTTP и WebSocket (Boost.Beast), а также другие средства, востребованные при построении высоконагруженых систем. Использование современных возможностей стандарта C++23 дополняет эти библиотеки: новые средства языка упрощают написание производительного кода и снижают вероятность ошибок. Отдельного внимания заслуживает шаблонное метапрограммирование, которое позволяет реализовать рефлексию на этапе компиляции и тем самым сократить нагрузку на этапе выполнения.
%
%На практике выбор конкретных решений не сводится к одной библиотеке или фреймворку. Существующие проекты, такие как Drogon и Pistache, демонстрируют разные подходы к построению серверных приложений на C++. Их анализ позволяет увидеть компромиссы между удобством разработки, производительностью и переносимостью. В работе рассматриваются эти решения и лежащие в их основе техники.
%
%Цель данного исследования — определить и оценить, в какой мере совмещение современных HTTP-технологий и C++, на примере библиотек Boost, стандарта C++23, приёмов шаблонного метапрограммирования и существующих фреймворков Drogon и Pistache, позволяет удобно разрабатывать высокопроизводительные серверные решения. Для этого предполагается проанализировать эргономику разработки на таком стеке и сравнить получаемые решения с существующими серверными фреймворками по задержке, пропускной способности, расходу памяти и трудозатратам разработчика.

\section{Библиотеки Boost как решение для \\кроссплатформенной интеграции и разработки \\высокопроизводительных систем}

\section{Использование современного стандарта языка С++23}

\section{Шаблонное метапрограммирование для рефлексии и сокращения нагрузки на этапе выполнения}

\section{Решения для разработки на С++: Drogon и Pistache}

\section{Экосистема Brazier}

\section{Заключение}

Brazier отличается от Drogon, Oat++ и Crow единой модульной архитектурой. Вместо разрозненных библиотек C++ он объединяет несколько ключевых компонентов. Внутри одной MVC-парадигмы работают асинхронная сеть на Boost.Asio, brazier-orm, реализующий трехзвенную ORM для PostgreSQL, brazier-logger, реализующий структурированное журналирование и обработка HTTP, HTTPS и WebSocket.

Brazier использует эргономичные абстракции C++20/23 чтобы уменьшить объем шаблонного кода. Они также убирают необходимость в ручном управлении необработанной памяти при маршрутизации и сопоставлении с базой данных. При этом безопасность типов сохраняется на этапе компиляции.

Brazier обеспечивает бесшовную работу с зависимостями и кроссплатформенный конвейер, используя инструментарий CMake и Vcpkg. Такая схема упрощает развертывание в Windows, Linux и macOS\@.

\pagebreak
\printbibliography{}

\end{document}